\subsection{\texorpdfstring{\(R^{n}\) 到拓扑群中的连续同态的局部定义}{R 的 n 次幂到拓扑群中的连续同态的局部定义}}

第 V 章 § 1 第 4 小节的命题 6 可以推广到所有的群 \(\mathbf{R}^n\) 。

命题 2. 设 A 是 \(\mathbf{R}^n\) 中包含 o 的一个超平行体；设 \(f\) 是 A 到拓扑群 G（写成乘法）中的连续映射，使得 \(f(x + y) = f(x)f(y)\) 对每对点 \(x, y\) （满足 \(x \in A\) 、 \(y \in A\) 与 \(x + y \in A\) ）成立。则存在 \(\mathbf{R}^n\) 到 G 中的唯一连续同态，它延拓 \(f\) 。

用第 V 章 § 1 第 4 小节命题 6 中的同样论证可以证明：延拓 \(f\) 的同态如果存在，则是唯一的。其次，子群 \(G_1\) （ \(G\) 的，由 \(f(A)\) 生成）是交换的；因为若 \(x\) 与 y 是 A 的任意两点，则 \(\frac{1}{2}x\) 、 \(\frac{1}{2}y\) 与 \(\frac{1}{2}(x+y)\) 都属于 A，因此

\[
f \left(\frac {1}{2} (\boldsymbol {x} + \boldsymbol {y})\right) = f \left(\frac {1}{2} \boldsymbol {x}\right) f \left(\frac {1}{2} \boldsymbol {y}\right) = f \left(\frac {1}{2} \boldsymbol {y}\right) f \left(\frac {1}{2} \boldsymbol {x}\right),
\]

这表明 \(f(\frac{1}{2} x)\) 与 \(f(\frac{1}{2} y)\) 交换；因此 \(f(x) = (f(\frac{1}{2} x))^2\) 与 \(f(y) = (f(\frac{1}{2} y))^2\) 也交换，这就证明了断言。

设 \(a_1, a_2, \ldots, a_n\) 是 \(n\) 个非零向量（含于 \(A\) 中且与超平行体的基向量成比例），且对每个指标 \(i\) ，设 \(D_i\) 是过 \(o\) 与 \(a_i\) 的直线，即诸点 \(ta_i\) （当 \(t\) 取遍 \(R\) 时）所成的集合。设 \(A_i\) 是一切 \(t \in R\) （使 \(ta_i \in A\) ）所成的集合；则 \(A_i\) 是 \(R\) 的一个包含 \([o, i]\) 的区间，而函数

\[
f _ {i} (t) = f \left(t \boldsymbol {a} _ {i}\right)
\]

在 \(A_{i}\) 上有定义且连续，并满足关系

\[
f _ {i} (t + t ^ {\prime}) = f _ {i} (t) f _ {i} (t ^ {\prime})
\]

只要 \(t, t'\) 与 \(t + t'\) 都属于 \(A_i\) 。由第 V 章 § 1 第 4 小节命题 6，存在一个连续同态 \(\overline{f}_i\) （ \(\mathbf{R}\) 到 \(G\) 中的），它延拓 \(f_i\) 。由于 \(\mathbf{R}^n\) 是诸子群 \(D_i\) 的直和，我们可以定义一个同态 \(\overline{f}\) （ \(\mathbf{R}^n\) 到交换群 \(G_1\) 中的），其规则为 \(\overline{f}(x) = \prod_{i=1}^{n} \overline{f}_i(t_i)\) （其中 \(x = \sum_{i=1}^{n} t_i a_i\) ）； \(\overline{f}\) 是 \(f\) 的延拓，因为若 \(x \in A\) ，则诸分量 \(x_i\) （ \(x\) 在诸直线 \(D_i\) 上的）也都属于 \(A\) （由诸 \(a_i\) 的取法）；此外， \(\overline{f}\) 在 \(R^n\) 上连续，因为它在每条直线 \(D_i\) 上连续，且 \(x_i\) 是 \(x\) 的线性函数（因而连续）。

推论 1. 设 V 是 \(R^{n}\) 中 o 的一个邻域，f 是 V 到拓扑群 G 中的连续映射，使得

\[
f (\boldsymbol {x} + \boldsymbol {y}) = f (\boldsymbol {x}) f (\boldsymbol {y})
\]

对每对点 \(x, y\) （满足 \(x \in V\) 、 \(y \in V\) 与 \(x + y \in V\) ）成立。则存在 \(\mathbf{R}^n\) 到 \(G\) 中的唯一连续同态，它与 \(f\) 在一个邻域 \(W\) （ \(o\) 的）的所有点处一致。

取 W 为以 o 为中心、含于 V 中的开长方体，并对 W 应用命题 2。

我们以后将看到， \(R^{n}\) 的这个性质可以推广到更大的一类拓扑群，即“单连通”群。

推论 2. 设 \(f\) 是 \(\mathbf{R}^n\) 与拓扑群 \(\mathbf{G}\) 之间的一个局部同构。则存在 \(\mathbf{R}^n\) 到 \(\mathbf{G}\) 的一个开子群上的唯一严格态射，它与 \(f\) 在 \(\mathbf{o}\) 的某个邻域的所有点处一致。

设 \(\overline{f}\) 是 \(\mathbf{R}^n\) 到 \(\mathbf{G}\) 中的、在 o 的某个邻域的所有点处与 \(f\) 一致的连续同态；由假设， \(\overline{f}(\mathbf{R}^n)\) 包含 \(\mathbf{G}\) 的单位元的一个邻域，因此（第 III 章 § 2 第 1 小节命题 4 的推论）是 \(\mathbf{G}\) 的一个开子群；此外，由第 III 章 § 2 第 8 小节命题 24， \(\overline{f}\) 是 \(\mathbf{R}^n\) 到 \(\overline{f}(\mathbf{R}^n)\) 上的一个严格态射。

定理 1. 每个局部同构于 \(\mathbf{R}^n\) 的连通群 G 都同构于某个群 \(\mathbf{R}^p \times \mathbf{T}^{n-p}\) （ \(0 \leqslant p \leqslant n\) ）。

适当选取的局部同构 \(f\) （在 \(\mathbf{R}^n\) 与 \(G\) 之间）可以延拓为 \(\mathbf{R}^n\) 到 \(G\) 的一个开子群上的严格态射（命题 2 的推论 2），并由于 \(G\) 连通而延拓到 \(G\) 自身。由此得出 \(G\) 同构于一个商群 \(\mathbf{R}^n / H\) （ \(\mathbf{R}^n\) 的）： \(H\) 是离散的，否则会存在点 \(x \neq 0\) （属于 \(H\) ，且任意接近 \(o\) ）使得 \(f(x) = f(0)\) ，这与 \(f\) 是局部同构这一假设相矛盾。因此定理是 § 1 第 1 小节定理 1 的推论。

